\subsection{李群情形}

命题 6. 设 k 是 R、C 或特征为 0 的非离散完备超度量域。设 G 是 k 上的一个有限维李群，e 是其单位元，g 是其李代数，h 是 g 的一个嘉当子代数，\(h_{r}\) 是 g 中属于 h 的正则元所成的集。

\begin{enumerate}
\def\labelenumi{(\roman{enumi})}
\item
  设 \(\mathfrak{s}\) 是 \(\mathfrak{h}\) 在 \(\mathfrak{g}\) 中的一个向量空间补，\(\mathfrak{s}_0\) 是 \(\mathfrak{s}\) 中 0 的一个在其上定义了指数映射的邻域，且 \(h_0 \in \mathfrak{h}_r\)。映射 \((s, h) \mapsto \mathrm{F}(s, h) = (\exp \operatorname{ad} s).h\)（从 \(\mathfrak{s}_0 \times \mathfrak{h}\) 到 \(\mathfrak{g}\)）在 \((0, h_0)\) 处是平展的。
\item
  映射 \((g,h)\mapsto \mathrm{F}'(g,h) = (\operatorname {Ad}g).h\)（从 \(\mathbf{G}\times \mathfrak{h}_r\) 到 \(\mathfrak{g}\)）是一个浸没。特别地，其像 \(\Omega\) 是开集。对所有 \(x\in \Omega\)，\(\mathfrak{g}^0 (x)\) 是 \(\mathfrak{g}\) 的一个与 \(\mathfrak{h}\) 在 \(\operatorname {Ad}(\mathrm{G})\) 下共轭的嘉当子代数
\item
  设 \(h_0 \in \mathfrak{h}_r\)。对任意邻域 \(\mathrm{U}\)（\(e\) 在 \(\mathrm{G}\) 中的一个邻域），集合 \(\bigcup_{a \in \mathrm{U}} (\mathrm{Ad}a)(\mathfrak{h}_r)\) 是 \(h_0\) 在 \(\mathfrak{g}\) 中的一个邻域。
\end{enumerate}

设 \(h_0\) 与 \(\mathfrak{s}\) 如 (i) 中所述。设 \(\mathrm{T}\) 是 \(\mathrm{F}\) 在 \((0, h_0)\) 处的切线性映射。于是 \(\mathrm{F}(0, h) = h\) 对所有 \(h \in \mathfrak{h}\) 成立，故 \(\mathrm{T}(0, h) = h\) 对所有 \(h \in \mathfrak{h}\) 成立。另一方面，当 \(\mathfrak{s}_0\) 充分小时，映射 \(s \mapsto \exp s\)（从 \(\mathfrak{s}_0\) 到 \(\operatorname{End}(\mathfrak{g})\)）在 0 处的切线性映射就是映射 \(s \mapsto \operatorname{ad}s\)（从 \(\mathfrak{s}\) 到 \(\operatorname{End}(\mathfrak{g})\)）。于是 \(\mathrm{T}(s, 0) = [s, h_0]\) 对所有 \(s \in \mathfrak{s}\) 成立。而从 \(\mathfrak{g}/\mathfrak{h}\) 到 \(\mathfrak{g}/\mathfrak{h}\) 的、由 \(\operatorname{ad}h_0\) 过渡到商所诱导的映射是双射的。由此推出 \(\mathrm{T}\) 是双射，从而得出 (i)。由于 \(\exp s = \operatorname{Ad}\exp s\) 对所有充分接近 0 的 \(s \in \mathfrak{s}\) 成立，(iii) 与 (ii) 的第一个断言随之成立。每个 \(x \in \Omega\) 都形如 \((\operatorname{Ad}a)(h)\)，其中 \(a \in G\)，\(h \in \mathfrak{h}_r\)，于是 \(\mathfrak{g}^0(x) = (\operatorname{Ad}a)(\mathfrak{g}^0(h)) = (\operatorname{Ad}a)(\mathfrak{h})\) 是 \(\mathfrak{g}\) 的一个与 \(\mathfrak{h}\) 在 \(\operatorname{Ad}(G)\) 下共轭的子代数。

